%%%%%%%%%%%%%%%%%%%%%%%%%%%%%%%%%%%%%%%%%
%
% CMPT 435
% Lab Zero
%
%%%%%%%%%%%%%%%%%%%%%%%%%%%%%%%%%%%%%%%%%

%%%%%%%%%%%%%%%%%%%%%%%%%%%%%%%%%%%%%%%%%
% Short Sectioned Assignment
% LaTeX Template
% Version 1.0 (5/5/12)
%
% This template has been downloaded from: http://www.LaTeXTemplates.com
% Original author: % Frits Wenneker (http://www.howtotex.com)
% License: CC BY-NC-SA 3.0 (http://creativecommons.org/licenses/by-nc-sa/3.0/)
% Modified by Alan G. Labouseur  - alan@labouseur.com, and Patrick Tyler - Patrick.Tyler1@marist.edu
%
%%%%%%%%%%%%%%%%%%%%%%%%%%%%%%%%%%%%%%%%%

%----------------------------------------------------------------------------------------
%	PACKAGES AND OTHER DOCUMENT CONFIGURATIONS
%----------------------------------------------------------------------------------------

\documentclass[letterpaper, 10pt]{article} 

\usepackage[english]{babel} % English language/hyphenation
\usepackage{graphicx}
\usepackage[lined,linesnumbered,commentsnumbered]{algorithm2e}
\usepackage{listings}
\usepackage{fancyhdr} % Custom headers and footers
\pagestyle{fancyplain} % Makes all pages in the document conform to the custom headers and footers
\usepackage{lastpage}
\usepackage{url}
\usepackage{xcolor}
\usepackage{titlesec}
\usepackage{xurl}
\usepackage{multicol}

% Stolen from https://www.overleaf.com/learn/latex/Code_listing 
\definecolor{codegreen}{rgb}{0,0.6,0}
\definecolor{codegray}{rgb}{0.5,0.5,0.5}
\definecolor{codepurple}{rgb}{0.58,0,0.82}
\definecolor{backcolour}{rgb}{0.95,0.95,0.92}

\lstdefinestyle{mystyle}{
    backgroundcolor=\color{backcolour},   
    commentstyle=\color{codegreen},
    keywordstyle=\color{magenta},
    numberstyle=\tiny\color{codegray},
    stringstyle=\color{codepurple},
    basicstyle=\ttfamily\footnotesize,
    breakatwhitespace=false,         
    breaklines=true,                 
    captionpos=b,                    
    keepspaces=true,                 
    numbers=left,                    
    numbersep=5pt,                  
    showspaces=false,                
    showstringspaces=false,
    showtabs=false,                  
    tabsize=2
}
\lstset{style=mystyle, language=c++}


\fancyhead{} % No page header - if you want one, create it in the same way as the footers below
\fancyfoot[L]{} % Empty left footer
\fancyfoot[R]{}

\renewcommand{\headrulewidth}{0pt} % Remove header underlines
\renewcommand{\footrulewidth}{0pt} % Remove footer underlines
\setlength{\headheight}{13.6pt} % Customize the height of the header

%----------------------------------------------------------------------------------------
%	TITLE SECTION
%----------------------------------------------------------------------------------------

\newcommand{\horrule}[1]{\rule{\linewidth}{#1}} % Create horizontal rule command with 1 argument of height

\title{	
   \normalfont \normalsize 
   \textsc{CMPT 308 - Spring 2024 - Dr. Labouseur} \\[10pt] % Header stuff.
   \horrule{0.5pt} \\[0.25cm] 	% Top horizontal rule
   \huge Lab 1 -- ~Data \& Information\\     	    % Assignment title
   \horrule{0.5pt} \\[0.25cm] 	% Bottom horizontal rule
}

\author{Patrick Tyler \\ \normalsize Patrick.Tyler1@marist.edu}

\date{\normalsize\today} 	% Today's date.

\begin{document}

\maketitle % Print the title

%----------------------------------------------------------------------------------------
%   CONTENT SECTION
%----------------------------------------------------------------------------------------

% - -- -  - -- -  - -- -  -
\section{Data vs Information}
Data is simply "stuff" in the world.
Data can take the form of anything... numbers, lines, or words.
Data, however, cannot make conclusions. 
Information is data empowered with the context to make useful statements
   the trends or lack there of.
For example, if we had the numbers 98 90 70 63 they are meaningless without context;
they carry no information, but they are still perfectly valid data.
Now let's add some context: those numbers are 4 grades from a professors class of 20.
Now this data means something, it represents some information.
We could make statements such as the average of these grades are 80.25 
   because it makes to average these numbers in the given context.
We could even try to make farther reaching statement such as the average
   for the entire class is probably around 80.25.
Still, we must be careful with how much information we actually have and
   whether can can acurately make certain assumptions.
In example, the context and information could change: let's say theses were
   the top 4 student grades.
The data remains the same but the context changes everything.
There a myriad of conclusions you might want to make about this data 
   but have you considered all the information and especially 
   the missing information?
What is the grading scale?
Are these even final grades?
What is rigor of the class?
\newline
\newline
Databases attempt to organize data in efficient ways whose structure,
    most importantly, aimes to convey all of the respective information.
Organizations who maintain a lot of data want to embed their "business rules"
    directly into how they store their data, and database are perfect for this.
For example, let's look at the Marist Banner system.
(I not have access to the actual database so this is inferred from
    the api their frontend website uses, an example section json entry 
    is below for reference)
Firstly, Banner stores this "id" attribute. From the outside this seems like
    a completely obscure number with no meaning.
And, it is, this number does not hold any meaning in the way numbers generally
   have meaning.
It is likely only there to give each section a unique identifier.
Without being coupled with the information of being an "id" this data is a
   completely useless number.
Sections also likely have a one-to-many reltionship with "faculty" and
   "meetingsFaculty".
This relational idea gives information: there is a group of "faculty" 
   and "meetingsFaculty" which which can be apart of a section.
Which, from what we know about courses at Marist, makes sense.
Every section has one to many teachers and some meeting times.
Databases organize data in order applying the business logic of the organization
   forming their data into information.
\section{Data Models}
Hierarchical databases store their data through a tree-like structure conveying
   information through relationships such as parent and child. 
Every child can only have a single parent and a parent can have an indefinite
   of children.
These contrainsts nicely represent data that can form perfect trees.
However it breaks down when any other relationship needs to be represented. 
For instance, let's go back to the Marist database imagine a simplistic version
   of storing sections, meetings, and their teachers.
In a hierarchical database sections and meetings are a perfect match.
At least simplistically, every section has 0 to many meetings and every
   meeting has exactly one parent.
However, representing their teachers is difficult even if we simplify that
   each section only has a single teacher.
Are they children of section?
Then we would lose the information that children of sections are meetings.
Are they children of meetings?
This could work but then we would be storing professor data for each meeting.
Whenever data is repeatibly stored it not only takes up more space but even more
   worryingly can lead to update, insertion, and deletion anomalies.
\newline
\newline
Maybe another type of data models can help us!
Network databases permit a child to have more than one parent.
This means we can have one instance of professor and all of their 
   meetings point to them.
We maintain the structure/ information that children of sections are meetings
   as well as children of meetings are professors without duplication.
This is great until we have a professor that does not teach any meetings.
Where should we put this professor?
There is no good answer in either of these databases.
This is one of the fundamental reasons why relational databases are
   largely thought as better.
\newline
\newline
There is another type of storage method known as XML.
I personally have to appreciate XML as I use HTMX, a small javascript library which
   uses html attributes to send AJAX requests to an endpoint.
I can represent all of my needs to send information from my frontend to my backend.
However, I think it is better suited as a data transfer structure rather than
   a full data storage and query system because it runs into the same problems
   as network database. 
XML does allow multiple parents and children but still does not have a good place
   for data that is not connected.


\section{Proof of Life}
 
\begin{figure}[h]
  \centering
  \includegraphics[width=0.8\textwidth]{./assets/pg-admin.png}
  \caption{PG admin running on my computer showing the cap database}
\end{figure}

\begin{figure}[h]
  \centering
  \includegraphics[width=0.9\textwidth]{assets/postgre.png}
  \caption{Postgre sql installed}
\end{figure}


\begin{figure}[h]
  \centering
  \includegraphics[width=0.6\textwidth]{assets/psql.png}
  \caption{psql cli tool working in powershell }
\end{figure}
 
\pagebreak
\begin{multicols}{2}
\lstinputlisting[
     caption={A Course's section from Banner's api},
     label={api}]{./1lab/api.json}
\end{multicols}
\end{document}
